\documentclass[12pt]{article}
\usepackage[margin=1in]{geometry}
\usepackage{enumitem}
\usepackage{titlesec}
\usepackage{parskip}
\usepackage{graphicx}
\usepackage{hyperref}
\usepackage{fontenc}

\title{\textbf{Cox 1987 RG Notes}\\
\large Cox, Gary W. \textit{The Efficient Secret: The Cabinet and the Development of Political Parties in Victorian England}.\\
Cambridge: Cambridge University Press, 1987.}
\author{}
\date{}

\begin{document}

\maketitle

\section{Main Idea}



%--------------------------------------------------
\section{General Notes}
%--------------------------------------------------

\begin{itemize}[leftmargin=*, itemsep=4pt]
    \item ``We seek to explain the evolution of two crucial behavioral aspects of the modern English polity: the marked regularity with which MPs vote with their parties, and the strong tendency of voters to vote for parties rather than individuals.... To show when and explain why party became the dominant influence on voting behavior in England -- both in the legislature and in the electorate.''
    \item Two primary independent variables: ``the extension of the suffrage and the centralization of legislative authority.''
    \item Rapid expansion of private bills (see below section) as the 19th century progressed, largely stemming initially from the complexity of regional investments in railroad services. A great deal of lobbying became the norm as voters sought to win the favor of MPs, but as this became broadly known, also led to significantly more public debate through procedural means
    \item Mid-century parliament (post first and second Reform Acts) featured a great deal of party independence: legislators stood to gain significant clout by voting their conscience rather than their party, and party discipline plummeted. Yet, it subsequently significant rebounded later in the 19th century
    \item Primary hypotheses:
    \begin{enumerate}
        \item ``Since the cost of becoming informed increases with the complexity and frequency of votes, nineteenth century MPs should have been more likely to rely on free information -- a chief element of which was the recommendations of their party leaders -- as motions pushed to a division became more complex and frequent."
        \item ``Since the benefit to acquiring information depends on the chances that the information will actually change the vote decision, any process whereby nineteenth century MPs became more confident that their party's line would be acceptable to them were they to investigate the matter thoroughly should have increased party voting. Party cohesion should have increased when the probability of agreement between the party and other relevant groups increased.''
    \end{enumerate}
    \item Argument: as the information available to voters greatly increased, and the number of voters for whom that information was relevant increased as well, it became an informational challenge for voters to hold specific MPs accountable; instead, voters turned towards the parties as broader symbols of accountability, to be punished and rewarded collectively. Following this transition, it became incentive compatible for individual MPs to cede power/authority to the collective party, as its success would in turn secure their own position(s).
\end{itemize}


\section{Important Specific Factual Claims}

\begin{itemize}[leftmargin=*, itemsep=4pt]
    \item Private bill: dealing with a specific policy issue regarding a firm, individual, constituency, etc. -- NOT designed to establish blanket policy guidelines, but rather to solve a particular problem within the MP's constituency (explicitly or implied)
    \item Core developments in 19th-century Britain:
    \begin{itemize}
        \item Reform Act of 1832 \& Reform Act of 1867 provided significant expansions of the franchise and provided greater equality of that franchise (previously -- and still following 1832 -- many small constituencies were heavily overrepresented in Parliament). Further expansion and ``clean-ups'' in the third Reform Act
        \item Dramatic increase in population met by an unexpected proportional increase in real wealth: industrial revolution allowed for greatly increased productivity and output per worker and in turn lowered consumer costs \& facilitated broader consumer markets. Shift in population from agrarian to industrial focus
        \item Rapid expansion of number of news outlets and the number of citizens exposed to the press: lowered/eliminated duties on advertisements/stamps/papers, steam-driven printing press in 1811, development of the telegraph in the late 1840s. Rapid diffusion of political information, particularly during elections
        \item Development/expansion of steam locomotive and developed railroad system, providing both significant economic growth (\& complexity) and interpersonal transportation; this in turn forced procedural reforms in parliament with the a delineation between ``public'' and ``private'' bills (think: who benefits?).
    \end{itemize}
    \item The ``efficient secret'' itself is the fusion of legislative and executive authority within the cabinet, rather than as separate authoritative bodies. Core question is then what facilitated this delegation of power from individual MPs to their parties and, subsequently, the respective cabinet.
\end{itemize}


\section{Connections}

\begin{itemize}
    \item Olson 1965: party formation as a collective action solution. Each MP could retain their own influence/power/oversight, but in doing so would greatly hinder the collective legislative process -- individuals needed to be incentivized to acquiesce these individual benefits for the good of the party as a whole. (Is this directly cited?)
    \item Huntington 1968: political development as a response to economic/social development, very cleanly reflected in Cox's account of parliament
    \item Crisp et al 2004: legislators deciding to focus on either local or national issues -- either providers of public goods or providers of private goods but to their constituents. Connected to the ``private bill'' whereby legislators can choose to focus on more specifically/local policy concerns
    \item Carey \& Shugart 1995 -- under which electoral systems would legislators be more likely to pursue a personal vote (and subsequently less likely to cede power to the party / cabinet)?
    \item Jones \& Hwang 2005, Katz \& Mair 1995 -- Cartel theory: Cartel theory claims that the party has very limited control over its members; Cox suggests that this has been largely overcome. A driving question here (and good to write about) -- is this because cartel theory is disproven, or is it because it is still incentive-compatible for MPs to cede authority to the central party/cabinet, such that cartel theory is circumvented in a new equilibrium outcome?
\end{itemize}


\section{My Critiques}

\begin{itemize}
    \item
\end{itemize}


\end{document}
